% Source PDF: source/普通高考/2020/2020全国2文(甘肃,青海,内蒙古,黑龙江,吉林,辽宁.宁夏,新疆,陕西,重庆).pdf, page 1.
% PDF embedded-bitmap attempt: pdfimages -list found no image objects; the flowchart is vector line art.
% Original comparison crop: tmp/redraw_flowchart_2020_national_paper_2_liberal/q7_source_exact.png,
% cropped from a no-grid 300 dpi page render after grid100/grid50 localization.
% Elements: rounded start/end terminals, input/output parallelograms, two update rectangles,
% one decision diamond, directed connectors, the left loop-back branch, and 是/否 labels.
% Node→directed-edge list: 开始→输入(k,a)→a=2a+1→k=k+1→a>10;
% a>10 的 是→输出 k→结束；a>10 的 否→左侧回边→a=2a+1.
% solid segments: all node borders and connectors; dashed segments: none.
% labels: every node is locally zero-indented and horizontally/vertically centered;
% 是 sits beside the downward true branch and 否 sits beside the left return branch.

\begin{tikzpicture}[
  x=1cm,
  y=0.91cm,
  >=Stealth,
  line cap=butt,
  line join=miter,
  every node/.style={font=\normalsize,anchor=center,align=center,inner sep=0pt,
    execute at begin node={\setlength{\parindent}{0pt}}},
  flow/.style={draw=black,line width=0.45pt,anchor=center,align=center,inner sep=0pt,
    execute at begin node={\setlength{\parindent}{0pt}}},
]
  % The source is a schematic flowchart; coordinates preserve topology and compact spacing.
  \node[flow,rounded corners=0.16cm,minimum width=1.05cm,minimum height=0.64cm]
    (start) at (2.64,8.25) {开始};
  % Input/output parallelograms are explicit polygons so their text bbox stays centered.
  \path[flow] (1.55,7.50) -- (3.73,7.50) -- (3.53,6.78) -- (1.33,6.78) -- cycle;
  \node at (2.54,7.14) {输入 \(k,a\)};
  \node[flow,rectangle,minimum width=2.00cm,minimum height=0.62cm]
    (updatea) at (2.64,5.53) {\(a=2a+1\)};
  \node[flow,rectangle,minimum width=2.00cm,minimum height=0.62cm]
    (updatek) at (2.64,4.28) {\(k=k+1\)};
  \node[flow,diamond,aspect=2.75,minimum width=2.45cm,minimum height=0.72cm]
    (check) at (2.64,3.03) {\(a>10\)};
  \path[flow] (1.90,1.99) -- (3.40,1.99) -- (3.20,1.30) -- (1.65,1.30) -- cycle;
  \node at (2.54,1.64) {输出 \(k\)};
  \node[flow,rounded corners=0.16cm,minimum width=1.05cm,minimum height=0.64cm]
    (finish) at (2.64,0.33) {结束};

  % Main path; the true branch goes directly to output and then end.
  \draw[->,flow] (start.south) -- (2.64,7.50);
  \draw[->,flow] (2.64,6.78) -- (updatea.north);
  \draw[->,flow] (updatea.south) -- (updatek.north);
  \draw[->,flow] (updatek.south) -- (check.north);
  \draw[->,flow] (check.south) -- node[pos=0.28,right=0.06cm] {是} (2.64,1.99);
  \draw[->,flow] (2.64,1.30) -- (finish.north);

  % The false branch has its own left return path and re-enters updatea.
  \coordinate (loop-left) at (0.45,3.03);
  % The return arrow enters the main connector above the update box, as in the source;
  % keeping this junction clear avoids a duplicate arrowhead at the update border.
  % Keep the return shaft at least 2.5 mm above updatea's top edge.
  \coordinate (loop-top) at (0.45,6.35);
  % The return arrow terminates at the visible junction on the main shaft;
  % the shaft's own arrow continues into updatea.
  \coordinate (loop-bend) at (2.10,6.35);
  \coordinate (loop-merge) at (2.64,6.35);
  \draw[flow] (check.west) -- (loop-left) -- (loop-top);
  \draw[->,flow] (loop-top) -- (loop-bend) -- (loop-merge);
  \node[anchor=east,inner sep=0pt] at (1.33,3.77) {否};
\end{tikzpicture}
